\documentclass[../main]{subfiles}
\begin{document}
\ifSubfilesClassLoaded{\setcounter{page}{40}}{}
\section{Коммунизм и управление общественной жизнью}
\label{sec:07_kommunizm_upravlenie}

Производство и распределение свободного времени в переходный период к коммунизму нуждаются в государстве. Однако, государство необходимо только постольку, поскольку производство свободного времени совпадает с производством и распределением вещей, нужных людям, и поскольку само это производство вещей ещё не осуществляется по-коммунистически. Само же свободное время, выступая, собственно, как свободное (время коммунизма), выходит из-под всякого внешнего по отношению к нему государственного управления.

В пространстве свободного времени государства не существует даже сейчас, при капитализме. Там, где имеет место свободная коммунистическая деятельность, управление осуществляется всеми членами свободной ассоциации людей. При увеличении доли коммунистической деятельности, государство отмирает.

Способствование производству свободного времени и непосредственно- коллективных форм свободной деятельности является актом уничтожения условий для существования классов в переходный период, и в той же мере актом самоуничтожения государства. Государство — это лишь аппарат обеспечения интересов господствующего класса, осуществляющий управление в интересах этого класса, а эта его основная функция отпадает по мере уничтожения классов, в том числе и господствующего.

Задача государства в переходный период, таким образом, сводится к обеспечению ограничения (уничтожения) товарного характера производства и заботе о создании условий для коммунистической свободной деятельности членов общества.

Первая из них возникает и перед капиталистическими государствами. Капитализм постоянно порождает войны, а логика войны требует системы прямого распределения продуктов, предназначенных для личного потребления. Государство, не справляющееся с этой задачей, обречено на уничтожение (или само по себе, как государство, или вместе с населением). Задача налаживания такой системы предполагает в ряде пунктов притеснение отдельных капиталистов в интересах капиталистов как класса, что не всегда может быть осуществлено в полной мере и с достаточной степенью эффективности.

На войне личное потребление — как продуктов питания, так и вещевого довольствия, — не может обеспечиваться в порядке капиталистического частного характера присвоения продуктов труда. Солдат, даже если он наёмник, во время ведения войны должен быть обеспечен своим работодателем, как минимум, оружием, пропитанием и обмундированием. В противном случае армия превращается в разбойничьи шайки. Но эти шайки не способны вести современную войну. Для ведения войны современная армия нуждается в высокотехнологичном оружии, средствах личного потребления, необходимых для воспроизводства живой личности солдата, а также средствах личной защиты. И даже если это всё обеспечивается путём прямого грабежа населения захваченных территорий, это не может быть личным делом отдельных солдат. Это всего лишь условие того, чтобы он мог выполнять свою работу. Но это условие отличается от условий наёмного труда, основанного на частном характере присвоения. Присвоение продуктов труда во всех своих существенных и значительных для социального воспроизводства индивида моментах в классических капиталистических отношениях купли-продажи рабочей силы вынесено за рамки производства, основанного на наёмном труде. Личное потребление определяется частным характером присвоения продуктов труда, опосредованным деньгами. И это, в свою очередь, гарантирует, что работник снова и снова будет вынужден приступать к труду, для обеспечения себя всем необходимым. На войне дело кардинальным образом отличается. Потребление в этом случае не может быть даже формально отдано на откуп самих индивидов.

Таким образом, современная война это, кроме всего прочего, конкуренция выработанных ещё при капитализме форм нетоварного производства — непосредственного производства человека. Причём здесь имеется в виду и армия как собственно-военная сила, так и трудовая армия, по-разному формирующаяся. Акцент смещается. Если для нормального хода капитализма на первый план выступает то, что человек должен делать своё дело, чтобы жить, то здесь, наоборот, человек должен жить, чтобы делать своё дело. Государство вынуждено брать на себя ответственность за это. Причём, денежно-товарные отношения становятся очевидным препятствием для выполнения этих функций.

Война делает невозможными старые, довоенные способы и формы деятельности людей. Государство в таких условиях объективно усиливается, берёт на себя функции распределения в процессе непосредственного производства человека как в армии, так и в тылу, занимается налаживанием и поддержанием хозяйственных связей. Но государство как аппарат подавления одного класса другим, точнее, подавления меньшинством большинства, в принципе не может решить эти задачи в полной мере, и уж тем более не ставит себе те задачи, которые не служат интересам правящего класса. Но именно такого рода задач (расходящихся с интересами правящего класса) во время войны у общества становится очень много. Они объективно встают перед людьми потому, что от их решения зависит само их физическое существование. Государство эти проблемы часто даже и не берётся решать, оставляя их на откуп разнообразнейшим негосударственным организациям и объединениям людей. Это вынуждает миллионы людей искать новые формы организации, которые бы брали на себя государственные функции. И эти органы общественного самоуправления, берущие на себя функции государства, строго говоря, в своей тенденции и являются, и не являются государством.

Наиболее важная и формообразующая из проблем, решение которых требует новых организационных форм в противовес капиталистическому государству, это предотвращение и прекращение войны, а также минимизация для общества потерь и разрушений, ей вызванных. Реализация этой важнейшей задачи как раз и противоречит интересам капитала, который еще не извлек из войны всех выгод, что мог извлечь. Подчеркнём, что тут может иметь исторический смысл не война сама по себе, не насилие само по себе, а сопротивление насилию реакционного класса, осуществляемому группировками этого класса во имя своих корыстных интересов, да еще и руками другого класса. В сопротивлении и формах сопротивления (хотя не только в них), а не в войне самой по себе, кроется важнейшая тенденция к самоорганизации на подступах к революции. Без такой самоорганизации,     выработанные      в    рамках     капитализма      способы непосредственного, не опосредованного вещами-деньгами производства человека и экономические механизмы их осуществления, не могут стать отправной точкой перехода общества на новые основания развития.

Эти, порожденные тенденцией к войне и самой войной (как противостояние войне) формы самоорганизации для управления общественными являются переходными и могут двигаться в трёх направлениях: 1) интеграция с государством; 2) вырождение-исчезновение (интеграция не всегда означает вырождение); 3) взятие на себя всех функций государства, превращение в новую государственную власть. Эти организационные формы, таким образом, объективно являются к тому же формами классовой борьбы там, где они вынуждены противостоять государству, взяв на себя функции тех или иных государственных органов, а значит, лишить полномочий аппарат политического господства правящего класса. И, если в определённых условиях это фактически облегчение для этих органов, так как сами они выполнять свои функции не могут, то дальнейшее движение по этому пути неизбежно наталкивается на противостояние государства, на попытки ограничить функции и влияние форм самоорганизации населения. Каким будет исход и характер этого противостояния, зависит от конкретных обстоятельств борьбы и действий борющихся сторон.

Если рассматривать историю ХХ века, то одной из самых эффективных форм самоорганизации, берущих на себя функции государства, были советы. Советы рабочих, солдатских и крестьянских депутатов одновременно являлись представительно-законодательными, распорядительными и контрольными органами. И хотя советы возникли еще в 1905 году, реальное их разворачивание в государственную силу, двоевластие, и уж тем более лозунг «Вся власть советам», а также конкретные шаги от лозунга к его реализации могли возникнуть только в условиях политического банкротства старого государства в процессе изнуряющей войны. После Февральской революции банкротство государства как аппарата, претендующего на то, чтобы организовывать жизнь общества, стало очевидным, так как не только царское правительство, но и Временное правительство было бессильно в этом отношении. В то же время поднимались и активно включались в дела как местного самоуправления, так и государственного управления советы. Создавалась ситуация двоевластия как в тылу, так и в армии. В связи с этим особенно важно учитывать возможность в будущем ситуации, которая имела место весной и в начале лета 1917 года. Тогда возможен был мирный путь передачи власти советам как шаг по пути в сторону социальной революции — созданию условий для социализма. И возможность эта существовала благодаря двум обстоятельствам: силе и общественной значимости самих советов, и вооружённости населения в результате войны. Лозунг «Вся власть советам» в то время был констатацией наличия такого мирного пути революции. И только после июльских событий 1917 года мирный путь стал уже невозможен. Но сам факт такой возможности для нас наиболее важен. Ведь если сложатся такие условия, очень важно их не пропустить, тем более с учётом развития современного вооружения мирный путь, вполне может оказаться вообще единственно-возможным для революции.

В современных условиях могут родиться и рождаются другие формы, далеко не всегда связанные с войной, но всегда с необходимостью решить важные общественные проблемы. Суть этих переходных форм везде одна – взятие населением непосредственно на себя функций государственной власти, тем самым упразднение её в привычном для капитализма смысле.

Теперь перейдём к рассмотрению управления общественной жизнью при коммунизме. Управление общественной жизнью вообще сводится к планированию и корректировке плана с учётом реальных изменений, волеизъявления (распоряжения), учёта и контроля над управляемым процессом, непосредственного вмешательства в случае необходимости.

Планирование при коммунизме – это планирование развития человека, который становится основной целью общественного производства, планирование производства новых отношений между людьми. Такое планирование совпадает с прогнозированием, в том смысле, что самый надежный прогноз, это тот, к осуществлению которого прилагаются планомерные усилия. Все эти функции общество должно научиться выполнять без стоящего над ним органа «управления», а на самом деле угнетения, именуемого государством.

По мере же становления новой коллективности все его функции, которые относятся к производству вещей, должны быть переданы вещам. Государство в переходный к коммунизму период не просто отмирает, оно отмирает только в том случае, если работает на своё отмирание, в той мере, в какой оно заботится о том, чтобы люди обходились без него. Все остатки несвободы, в том числе и способы утилизации времени, с которыми связано существование государства, противостоят свободному времени как пространству человеческого развития, противостоят свободной человеческой деятельности. Там, где осуществима полная свобода, нет места для государства. Но там, где полная свобода ещё невозможна непосредственно, наличие определенного типа государства не только фиксирует факт этой невозможности, но и является инструментом для её преодоления.

Против отмирания государства выдвигается тот же аргумент, что и против снятия разделения труда. Суть его состоит в том, что общественная жизнь становится очень сложной и всё время усложняется, а потому для управления такой сложной жизнью нужна дифференциация функций и специальный аппарат. Но из этой сложности вовсе не следует, необходимость в особых, над обществом поставленных, отрядах чиновников и вооружённых людей. Не будь раскола общества на классы, самодействующая и самоуправляемая ассоциация людей была бы возможна, при всей её сложности, соответствующей уровню сложности общественной жизни, на чём вслед за Марксом и Энгельсом настаивал Ленин.

Тем более, что говоря о коммунизме, мы имеем в виду общество, где преодолено разделение труда, где люди более менее одинаково развиты в отношении здоровья, чувственности, мышления и нравственности, имеют политехническое, научно-теоретическое, универсальное образование на современном необходимом (то есть наивысшем) для общества уровне, а потому в короткие сроки могут осваивать любую частную «профессию» и прикладывать свои силы там, где это нужно, и где это им самим интересно, пусть даже один по преимуществу всю жизнь занимается изучением птиц, а другой балетом.

Но именно поэтому в управлении общественным развитием с использованием достигнутых технологий, опираясь на выработанную человечеством культуру мышления, при коммунизме участвуют все.              Это важнейшее условие универсального развития всех членов общества - на протяжении всей жизни все должны уметь управлять и иметь привычку управлять общественными делами – всеми делами, в которых сами принимают участие. Такое общество более реально,   чем    фантастический     классовый     мир     и   «гармонические взаимоотношения» между классами, пропагандируемые современными буржуазными идеологами. Наоборот, марксизм настаивает на том, что само существование государства говорит о том, что этот мир в условиях существования классов невозможен.

Для отмирания государства, должны быть соответствующие социальные и технические предпосылки. Что касается технических предпосылок, то с хорошим общеобязательным политехническим образованием на основе современной техники организовать управление общественными процессами было бы не так уж сложно в обществе, где отсутствует частная собственность (потому важнейшей предпосылкой является формальное обобществление средств производства). Конечно, для этого должны быть сняты общественные барьеры, что тоже требует от общества серьёзной работы и достижимо не сразу. Главное, что чисто технически уже на данном этапе это не представляется невозможным. Вся информация о состоянии производства и распределения во всех сферах общественной жизни в реальном времени, с возможностью прослеживания цепочек до самого начала, учётом ресурсов, потребностей технологических возможностей, возможностями связи и голосований, и т.д., и т.п. от низшего уровня (прямого вмешательства) до наивысшего коллективного принятия решений на уровне как локальных сообществ, так и всего общества, - всё в реальном времени может быть доступно каждому члену ассоциации.

Голосования и другие, привычные нам проявления демократии нужны только на первых порах. При переходе от управления людьми к управлению вещами и процессами необходимость в голосованиях понемногу отпадает, уступая место объективной необходимости, заключённой в вещах и во взаимодействии между ними. На этот момент указывал Ленин в «Государстве и Революции» подробно разбирая наследие Маркса и Энгельса: «Об "отмирании" и даже еще рельефнее и красочнее - о "засыпании" Энгельс говорит совершенно ясно и определенно по отношению к эпохе после "взятия средств производства во владение государством от имени всего общества", т.е. после социалистической революции. Мы все знаем, что политической формой "государства" в это время является самая полная демократия. Но никому из оппортунистов, бесстыдно искажающих марксизм, не приходит в голову, что речь идет здесь, следовательно, у Энгельса, о "засыпании" и "отмирании" демократии. Это кажется на первый взгляд очень странным. Но "непонятно" это только для того, кто не вдумался, что демократия есть тоже государство и что, следовательно, демократия тоже исчезнет, когда исчезнет государство. Буржуазное государство может "уничтожить" только революция. Государство вообще, т. е. самая полная демократия, может только "отмереть".

В современную эпоху бурного развития технологий, организация отмирания государства не представляется таким уж сложным делом чисто с технической стороны. Учёт и контроль в деле управления вещами и процессами в таких условиях становится делом техники, а его результаты доступными каждому. В этом деле отпадает необходимость иерархии и исчезает предпосылка для того, чтобы эти функции могли выделиться функциям в отдельную профессию. При внедрении автоматических систем для управления при обеспечении полной доступности к базам данных для каждого и работы с ними исчезает необходимость в «особой армии чиновников», которые могли бы даже «по- видимости» стоять «над», обществом как особый «аппарат».

Таким образом, будет решена одна из важнейших задач перехода к социализму – налаживание учёта и контроля, на важности которой постоянно настаивал Ленин «Учёт и контроль - вот главное, что требуется для «налажения», для правильного функционирования первой фазы коммунистического общества». Речь идёт о таком учёте и контроле над общественным производством, организация которого уже сама по себе означает организацию новых форм коллективности «Если действительно все участвуют в управлении государством, тут уже капитализму не удержаться. И развитие капитализма, в свою очередь, создаёт предпосылки для того, чтобы действительно все могли участвовать в управлении государством. К таким предпосылкам принадлежит поголовная грамотность, осуществлённая уже рядом наиболее передовых капиталистических стран, затем «обучение и дисциплинирование» миллионов рабочих крупным, сложным, обобществлённым аппаратом почты, железных дорог, крупных фабрик, крупной торговли, банкового дела и т. д. и т. п.»

Мы же можем ещё раз акцентировать, что к таким предпосылкам относится уже теперь развитие компьютерной техники, и заложенные в ней возможности, которые в принципе не могут быть использованы при капитализме, когда нет единого общественного хозяйства, а связь между производителями осуществляется через рынок. Эти возможности в организации управления общественной жизнью исследовал в своих разработках Виктор Михайлович Глушков, разрабатывающий Общегосударственную систему управления экономикой в СССР.      Глушков рассматривал управление общественным производством в целом, при помощи управляющих систем, где каждый имеет доступ к подробнейшей картине производства в реальном времени, а там, где человеческие решения не нужны, система действует автоматически по заранее заданным людьми схемам в соответствии с общественными целями.

Чем больше будут развиваться технологии, тем проще будет организовать управление и обеспечить участие в нём всех членов общества. Что касается содержательной стороны, то она требует соответствующего развития членов общества. Сложность состоит в том, что начинать движение к коммунизму нужно не с людьми коммунистического завтра, а с людьми, сформировавшимися в условиях разделения труда, преодолевающими свою ограниченность посредством образования и коллективной деятельности.

Политехническое образование становится необходимостью для членов свободной ассоциации. Идея политехнизма довольно хорошо разработана, так что мы вполне можем отослать читателя к соответствующей литературе. Здесь же отметим только, что для управления общественными процессами члены общества должны обладать целостной научной картиной мира и универсальным умением применения своих знаний. Современному чиновнику для выполнения своей функции научная картина мира даже в первом приближении совершенно не нужна, поэтому она у него сама собой исчезает, как что-то излишнее, даже если когда-то была. Чиновник - всего лишь живая частичная деталь устаревшей государственной машины. «Общество, которое по-новому организует производство на основе свободной и равной ассоциации производителей, отправит всю государственную машину туда, где ей будет тогда настоящее место: в музей древностей, рядом с прялкой и с бронзовым топором", - говорил Энгельс.

Конечно сама задача по управлению общественными процессами в такой ассоциации, и по сознательному построению общественной жизни будет требовать от людей очень многого, в том числе развитого диалектического мышления, развитой чувственности, нравственности, соответствующего уровня грамотности. Но это не недосягаемые высоты, доступные только для избранных, а то, что в соответствующих условиях станет достоянием каждого. По мере включения людей в активную общественную жизнь, как субъектов этой жизни, они будут выходить на передний план человеческой культуры. Так мы понимаем знаменитый афоризм Ленина о том, что каждая кухарка должна научиться управлять государством. Препятствием для этого выступает только способ деятельности современных людей, который превращает одного человека в государственного чиновника, а другого в кухарку.

Коммунистическое общество это - свободная ассоциация                 людей, управляющая своим собственным развитием по разумному плану.

Как именно она будет самоуправляться, с помощью каких именно средств и организационных мероприятий? На этот вопрос мы ответить не можем. Это - дело этой ассоциации, которая будет вынуждена с необходимостью развивать соответствующие организационные формы в конкретных исторических условиях. Самоуправление будет развиваться вместе с развитием коммунистического общества.
\end{document}
